\section{引言}

国际航运业正同时面临温室气体减排、燃料多元化和运行可靠性保障等多重压力。
IMO 2023 战略及行业订单趋势表明，替代燃料推进系统已由远期概念逐步转化为工程现实
\cite{imo2023ghgstrategy,classnk2025alternativefuels}。
在这一转型过程中，低压天然气、甲醇和氨等燃料被用于降低碳排放强度，或作为未来零碳燃料路线的重要候选方案。
因此，船用双燃料和多燃料发动机已成为大型动力系统研发中的重要方向
\cite{wingd_lng_faq_2025,wingd_ammonia_faq_2025,classnk_methanol_2024}。

然而，替代燃料的引入也带来了新的点火问题。
许多低反应性主燃料在实际船用发动机中仍需依赖可靠的高反应性点火源，
否则燃烧稳定性、循环重复性和负荷响应均可能受到影响。
换言之，燃料系统可以因替代燃料而具有更低碳、更灵活的特征，
但燃烧过程仍依赖一个尺度较小却十分关键的局部事件：
柴油预喷（pilot diesel injection，以下简称预喷）。

因此，柴油预喷仍是多燃料船用发动机中的关键使能机制，
其作用是在主燃料难以自行压燃的条件下提供稳定点火源
\cite{wingd_xdf_manual_2021,classnk_methanol_2024}。
在传统柴油发动机中，柴油本身是主要能量载体，压燃过程也已形成成熟的工程体系。
相比之下，在多燃料发动机中，预喷不再承担主燃料角色，而是成为适配替代燃料的关键点火接口。
对于喷油器开发而言，核心问题首先是如何控制预喷的空间发展：
pilot 贯穿距离、锥角和液相包络决定点火核在气缸内形成的位置、尺度和与主燃料充量的相对关系，
进而影响主燃料混合、火焰传播和燃烧效率
\cite{Karim2015DualFuel,WeiGeng2016DualFuelReview,Sahoo2009DualFuelReview}。
若 pilot 贯穿距离过短，点火核可能局限在不利区域或体积不足；
若贯穿距离过长，点火核与主燃料充量的相对位置也可能偏离设计意图，导致主燃料燃烧不充分，并增加液相撞壁的可能性。
因此，本文关注的贯穿距离预测并不只是为了避免失效，而是为了给预喷控制提供可快速查询的几何依据。

在这一控制目标未被满足时，预喷才会进一步表现为潜在失效来源。
若液相包络在形成有效点火核之前已触及活塞顶或缸壁，
则可能导致壁面润湿、润滑油稀释、局部混合恶化以及未燃碳氢排放升高等不利后果
\cite{peraza2022swi,ahmad2025fueldilution,MorieiraMotaPanao2010}。
在船用大缸径尺度下，该风险同样不容忽视。
喷嘴至壁面的几何距离、喷射参数、环境条件以及燃烧室尺度共同决定液相是否会在非预期区域持续存在。
因此，所需回答的问题不仅是预喷能否可靠点火，也不仅是是否撞壁，而是其空间发展是否落在有利于点火、混合和安全裕度的可控窗口内。

已有研究和工程实践已形成若干用于评估柴油撞壁风险的技术路线。
发动机试验最接近真实运行状态，能够在同一系统中考察喷雾、燃烧、壁面相互作用和排放结果。
高保真 CFD 可提供更丰富的流动细节，适用于研究喷雾、蒸发、混合和燃烧之间的耦合机制。
柴油喷雾研究也长期使用经验贯穿距离关联式和现象学模型快速描述喷雾发展，
其中贯穿距离与壁距的比较构成了一阶撞壁筛查的自然几何判据
\cite{HiroyasuArai1990,NaberSiebers1996,Baumgarten2006MixFormation}。
与此同时，光学喷雾实验能够在受控环境下提供可重复、可横向比较的液相包络证据，
高速 Mie 成像尤其适合观察液相边界和宏观贯穿距离演化 \cite{Linne2013SprayImaging}。

然而，上述工具各有适用边界，难以单独覆盖早期设计筛查的全部需求。
发动机试验成本高、周期长，且难以隔离单一喷嘴或喷射变量的影响。
高保真 CFD 信息丰富，但依赖复杂输入、模型调参与显著计算资源；
已有反应柴油喷雾 LES 研究显示，单次较高保真仿真即可消耗数万
CPU core-hours \cite{CurtisNiemeyerSung2017,liu2020sprayreview}。
经验模型结构紧凑且便于解释，但若要同时覆盖喷嘴开启/关闭瞬态、多孔喷嘴束间差异以及有限观测窗口，
其函数形式往往需要引入更多经验参数、分段规则或工况特定拟合，从而削弱快速泛化能力。
光学实验能够提供图像证据，但逐次后处理结果本质上仍为制表化离散观测；
若仅停留于观测结果本身，尚不能直接转化为设计迭代中可快速查询的条件参数化模型。

短预喷工况进一步放大了上述矛盾。
经典经验贯穿距离定律在物理上仍具有重要价值，但当观测窗口主要由喷射起始与喷射终止瞬态占据时，
其适用性会明显下降，或需以显著增加模型复杂度为代价 \cite{HiroyasuArai1990,ZhouLiYi2021,ZhouLiLaiWei2019}。
喷雾成像文献也表明，贯穿距离、锥角和边界位置并非脱离方法而存在的绝对真值，
而是由诊断方式、阈值规则和后处理定义共同生成的操作性观测量，
即依赖具体测量协议定义的可重复量，而非独立于方法的物理真值
\cite{MagnottiGenzale2013,JohnsonNaberLee2012,ecn_liquid_penetration,ecn_liquid_penetration_accessory}。
因此，Mie 视频的价值并不在于其能够天然给出完整液相物理场，
而在于其能够在给定光学布置和后处理协议下，稳定提供与工程筛查相关的几何线索，即可观测喷雾边界及其衍生指标。

然而，将这些几何线索转化为代理模型训练目标，还必须处理若干具体困难。
Mie 原始强度信号会受到照明、眩光、多重散射和后处理阈值的影响，导致不同喷油束的光学成像条件并不均匀；
多孔喷嘴使图像中同时存在多个方向相近但需要分别后处理的喷油束。
固定且有限的视场会使部分晚期贯穿距离观测转化为右删失样本；
若直接将原始轨迹的可见末端作为真实终点，模型便可能将``离开视场''误学为``物理饱和''。

由此可见，本文所需的并非单纯增加图像数据以提高数据驱动模型精度，
也并非直接训练更复杂的端到端图像到几何量黑箱回归器。
更关键的是建立一条可审计的工作流，将视频数据、几何定义、删失处理与不确定性表达有机连接起来。

数据驱动代理模型为这一目标提供了可行路径。
近期研究表明，神经网络可以基于工况参数预测喷雾贯穿距离和锥角等宏观特性，
从而在完整 CFD 与经验模型之间形成快速预测层 \cite{VazKarathanasis2026,liu2020sprayreview}。
对于本文的数据条件而言，大规模 Mie 视频还提供了额外机会：
同一光学配置下积累的多批次记录可被转化为统一的逐喷油束时序观测量，
并用于训练从工况条件到贯穿距离轨迹的映射。

然而，现有文献对本文所需的特定技术组合仍覆盖不足：
短脉冲预喷轨迹、有限视场导致的右删失、多孔喷嘴喷油束校准，
以及可供下游几何风险筛查直接使用的不确定性感知输出，尚未在同一框架中得到系统处理。
换言之，当前缺少的并非喷雾研究本身，也并非某一种单点测量技术，
而是位于光学实验归档与高成本验证之间、可快速查询、可重复执行且边界明确的中间层预测器。

因此，本文旨在评估并构建一条可审计的数据工作流：
将大规模 Mie 散射喷雾视频归档转化为位于光学实验归档与高成本验证之间的早期工程预测层，
从而为后续 CFD 或发动机试验快速筛选候选参数。
为实现这一目标，论文需要完成四项任务：
\begin{enumerate}
  \item 从高速非反应、非蒸发 Mie 视频中提取逐喷油束几何观测量；
  \item 构造能够缓解有限视场带来的右删失影响的贯穿距离轨迹目标；
  \item 训练能够同时输出贯穿距离均值与条件不确定性的代理模型；
  \item 将预测分布映射到二维简化发动机气缸几何中，形成概率性撞壁筛查代理指标。
\end{enumerate}

\subsection{主要贡献}

本文的主要贡献如下：
\begin{enumerate}
  \item 提出一套 CUDA 加速的大批量图像处理工作流，将 TB 级 Mie 散射视频转换为逐喷油束的时序观测量和边界轨迹；
  \item 提出一套降阶拟合策略，用于轨迹对齐、稳健外推以及缓解贯穿距离目标中的右删失影响；
  \item 提出一套多阶段不确定性感知学习框架，
  将确定性 MSE 训练、异方差 NLL 训练与基于可靠观测原始轨迹的教师引导细化结合起来，
  并最终得到在外部清洗协议下四次重复训练平均 RMSE 为~\SI{6.094}{mm} 且具有保守不确定度覆盖的 anchor-off Stage~3 代理模型；
  \item 提出一种简化的概率性撞壁筛查方法，在明确原型边界的前提下，将贯穿距离预测与下游几何碰撞评估解耦。
\end{enumerate}

本文的研究范围经有意收窄。
论文关注受控加压舱条件下非反应、非蒸发的柴油预喷喷雾，
聚焦于光学布置所观测到的宏观液相包络几何及其时间演化，而非完整的缸内燃烧过程。

因此，燃烧化学、蒸发、双燃料点火耦合、真实活塞顶与缸壁的热相互作用，
以及真实发动机壁膜真值预测均不在本文覆盖范围内。本文提出的方法应被理解为一种早期工程筛查代理模型，
而非已经完成验证的发动机撞壁真值模型。

此外需要明确指出，基于 MLP 的代理模型主要面向训练数据覆盖范围内的插值。
尽管本文对部分输入特征和工况编码进行了工程优化，但对于真正分布外的工况，预测均值与不确定度估计均不具备同等级别的可信度。
因此，该模型更适合被理解为数据域内的工程筛查工具，而非任意未知工况下均可信赖的通用喷雾预测器。

\subsection{论文结构}

本文余下部分结构如下。
第 2 章回顾文献并建立光学观测量、目标构造选择与分阶段学习策略的技术依据。
第 3 章描述实验数据与图像处理工作流，涵盖 OSCC 数据、人工标定、CUDA 预处理、喷油束配准、条带提取、CDF 贯穿距离和边界度量。
第 4 章描述轨迹目标、代理模型建模与概率性筛查，涵盖轨迹清洗、降阶拟合、特征设计、Stage~1--3 学习流程、基线构造和简化撞壁代理指标。
第 5 章报告实验结果，包括原始留出切分、外部清洗协议、Gaussian process 与物理基线、CRPS 和留一喷嘴诊断、试验批次层面的差异以及消融证据。
第 6 章讨论最终 Stage~3 模型为何是最可信的筛查模型、哪些内容已经得到验证，以及当前适用边界止于何处。
第 7 章给出结论，将已验证结果、当前原型边界与下一步开发重点分开陈述。

